% Hand-written narrative. Numbers here are transcribed from paper/results.json,
% paper/ablation_credit.tex and paper/ablation_arch.tex, all of which are
% regenerated by scripts/. Prose that argues a direction is written by hand on
% purpose: the direction is what the experiments decide, and this study reversed
% it three times before the measurements were trustworthy.

\newcommand{\ResultsSummarySentence}{%
Combining the closed-form per-agent credit with log-energy shaping is
significantly better than either alone ($d=-2.70$, $p=0.002$ and $d=-1.97$,
$p=0.009$ over six seeds), but only once the exploration scale is matched to the
actuation amplitude that actually damps the wing; at the conventional scale the
same effect is invisible, and the elaborate communication and action machinery we
derived from the same physics does not help at all.}

\newcommand{\ResultsNarrative}{%
Table~\ref{tab:main} and Figure~\ref{fig:comparison} give the full grid. Three
things in it matter.

First, the classical ladder separates cleanly by information structure rather
than by sophistication. Centralised LQG reaches $-5.52$ per second on healthy
conditions; the fully decentralised LQG, which sees exactly what the agents see,
reaches $-2.05$. Decentralisation costs a factor of $2.7$ on this plant. That
factor is a property of the problem, not of any learning algorithm, and it is why
the decentralised controller rather than the centralised one is the reference for
a distributed policy.

Second, the best learned policy reaches $-1.53$ per second, about three quarters
of the decentralised classical result. The average conceals a more interesting
pattern. Broken out by flight condition the decentralised LQG achieves $-5.42$,
$-1.07$ and $+0.32$ at $1.05$, $1.25$ and $1.45\,\Uf$ -- it fails outright at the
top of the envelope -- against $-1.98$, $-1.92$ and $-0.69$ for the learned
policy. The learned controller is far more uniform across the envelope and holds
a condition where the classical one diverges, while losing badly where the
classical one is strongest.

Third, and least comfortable: the robustness advantage we measured at a
conventional exploration scale does not survive at the scale that maximises
damping. At $\sigma=0.30$ the learned policies held six of the six feasible
conditions, matching a fault-aware oracle given the full state and the identity
of the failed surface. At $\sigma=0.05$ they hold five, the same as the classical
controllers, while damping three times better. Exploration scale trades
robustness against nominal performance on this task: broad exploration produces
policies that cope with failures they were not told about, narrow exploration
produces policies that damp precisely. We report both rather than the more
flattering one.}

\newcommand{\ResultsAblations}{%
\paragraph{Credit assignment.}
Six seeds per variant at $\sigma=0.05$ (Table~\ref{tab:ablation-credit}). The
per-agent credit alone reaches $-0.566 \pm 0.123$ and the shared log-energy
shaping alone $-0.690 \pm 0.145$; the two are statistically indistinguishable
($p=0.14$). Their combination reaches $-1.085 \pm 0.243$, significantly better
than either ($d=-2.70$, $p=0.002$ against credit alone; $d=-1.97$, $p=0.009$
against shaping alone). This is the exact-plus-residual structure the derivation
predicts: the closed form supplies the part of the credit the physics determines
exactly, and the shaping supplies the long-horizon part that an instantaneous
power balance cannot see. Neither is sufficient by itself.

\paragraph{Architecture.}
The mechanisms we derived from the same physics do not survive contact with
measurement (Table~\ref{tab:ablation-arch}). A plain feedforward policy with the
blended credit reaches $-1.085$; adding a recurrent phasor encoder gives $-0.487$,
adding phasor consensus gives $-0.279$, adding the phase-locked action head gives
$-0.561$, and sharing raw neighbour observations gives $-0.583$. Two of these
differences are significant against the plain policy ($d=-4.04$, $p=0.0006$ for
consensus; $d=-2.62$, $p=0.006$ for the full architecture). Every elaboration we
proposed made things worse, and we report this as the result it is.

We do not think this settles the question of whether communication helps
distributed flutter suppression in general. Three surfaces on a six-metre wing,
with air data already broadcast, is a small coordination problem, and at eight
surfaces the picture changes qualitatively -- every condition becomes feasible
and the learned policies hold all nine. It does settle that on this problem the
credit signal was worth deriving from the physics and the message content was
not.

\paragraph{Exploration scale.}
Figure~\ref{fig:exploration} is the finding with the widest reach. Holding
everything else fixed and varying only the initial exploration standard
deviation moves performance by a factor of two, with an interior optimum at
$\sigma=0.05$ and a sharp transition above it. The mechanism is a scale mismatch
that can be read off the axis: a controller achieving the classical result uses
about $0.1$ degrees RMS of deflection, while $\sigma=0.30$ corresponds to $4.5$
degrees. The policy's own exploration is then some forty times the amplitude that
works, and no credit decomposition can recover a signal buried underneath it.

This is not a tuning remark. Before we found it, every variant plateaued between
$-0.03$ and $-0.82$ regardless of credit signal, communication, action
parameterisation, memory, auxiliary supervision or base controller, and the
variant-to-variant differences that motivated several rounds of algorithmic work
were of the same order as the seed-to-seed spread. A capability probe rules out
the competing explanation: training on the single easiest condition gives $+0.21$
and $+0.27$, worse than training on the full distribution, so the plateau was
never a matter of the task being too broad.}

\newcommand{\PhaseNarrative}{%
The instrument separates the classical ladder sharply, which is what makes it
usable. Collocated feedback is essentially flat and responds at $5.50$ hertz,
never engaging the unstable mode at all -- the mechanism behind its poor
performance rather than a separate fact about it. Centralised LQG develops about
$90$ degrees of phase across the semispan and operates at $11.00$ hertz, against
the $11.03$ hertz the eigenvalue analysis predicts.

The same measurement also explains why authority is the wrong thing to look at.
Local rate feedback uses $44.7$ per cent of available deflection to achieve
$-0.03$ per second; centralised LQG uses $0.7$ per cent to achieve $-5.52$. The
better the controller, the less authority it spends. Flutter suppression is a
phase-precision problem, and the learned policies are not short of authority --
they apply a comparable amount at the wrong phase.}

\newcommand{\ConclusionSentence}{%
The measured lesson is narrower than the design we started from and more useful.
The credit decomposition was worth deriving from the physics and is significantly
better than either of its parts; the communication and action-parameterisation
mechanisms derived from the same physics were not, and a plain policy beat all of
them. Both conclusions were invisible until the exploration scale was matched to
the actuation amplitude that damps the wing, which on this problem mattered more
than any algorithmic choice we made.}
